\chapter{The permitted library base}\label{ch:base}
\module{(none --- this chapter is a contract, not mathematics)}

\section{The pin}

The permitted base is Lean \texttt{v4.34.0-rc2} and mathlib commit
\code{85e3a25e006c35636f0e53b0e9296caca2685bc0}, as recorded in the
repository's \texttt{lean-toolchain} and \texttt{lake-manifest.json}.
Everything in this document that is not on the following list is derived
here.

\section{What may be assumed}

\begin{enumerate}[label=(B\arabic*),leftmargin=*]
\item\label{base:linalg} Real and finite-dimensional linear algebra,
  Euclidean topology, compactness, smooth partitions of unity
  (\texttt{SmoothPartitionOfUnity}, \texttt{SmoothBumpCovering}), and
  elementary convexity.
\item\label{base:lp} Completeness of $L^p$ spaces; weak sequential
  compactness of bounded sequences in a Hilbert space.
\item\label{base:measure} Lebesgue and Hausdorff measure, Bochner
  integration, Fubini--Tonelli, monotone and dominated convergence,
  Radon--Nikodym for \emph{positive} measures, the \emph{positive-functional}
  Riesz--Markov theorem (\texttt{RealRMK.rieszMeasure} together with
  \code{RealRMK.integral\_rieszMeasure}), and regularity of the resulting
  positive measures.
\item\label{base:alaoglu} Banach--Alaoglu (\code{WeakDual.isCompact\_closedBall})
  and its sequential form for duals of separable normed spaces
  (\code{WeakDual.isSeqCompact\_closedBall}).
\item\label{base:besicovitch} The Besicovitch covering and differentiation
  theorems for Euclidean balls, including differentiation of one Radon
  measure with respect to another
  (\code{Besicovitch.exists\_disjoint\_closedBall\_covering\_ae},
  \code{Besicovitch.ae\_tendsto\_rnDeriv}).
\item\label{base:calculus} Finite-dimensional Fr\'echet calculus, the inverse
  and implicit function theorems
  (\code{HasStrictFDerivAt.to\_localInverse}), Taylor's theorem,
  Rademacher's theorem (\code{LipschitzWith.ae\_differentiableAt}), and
  \emph{equidimensional} Jacobian change of variables
  (\code{integral\_image\_eq\_integral\_abs\_det\_fderiv\_smul}).
\item\label{base:ode} Picard--Lindel\"of local existence and uniqueness,
  including a locally simultaneous family of solutions Lipschitz in the
  initial point
  (\code{IsPicardLindelof.exists\_forall\_mem\_closedBall\_eq\_hasDerivWithinAt\_lipschitzOnWith}),
  together with \texttt{ODE\_solution\_unique}.
\item\label{base:fourier} The Fourier transform on Schwartz functions,
  Fourier inversion, differentiation under the transform, and the $L^2$
  Plancherel isometry.  The pinned APIs include
  \code{MeasureTheory.Integrable.fourierInv\_fourier\_eq}, the derivative
  identities in \code{Analysis/Fourier/FourierTransformDeriv.lean}, and
  the $L^2$ Fourier linear isometry constructed in
  \code{Analysis/Fourier/LpSpace.lean}, with norm theorem
  \code{MeasureTheory.Lp.norm\_fourier\_eq}.  Real-valued Sobolev functions
  are sent to the complex theory by the canonical isometric inclusion.
  The inversion theorem requires integrability of the function and its
  Fourier transform, and continuity at the evaluation point.  In the
  Sobolev application it is first applied to smooth compactly supported
  approximants, for which these hypotheses hold, and the result is then
  extended by the quantitative Sobolev estimate and completeness.
\end{enumerate}

\section{What may \emph{not} be assumed}

The following are \emph{absent} from the base and are derived in this
document.  Each entry names the chapter that supplies it.  The list exists
because each of these has previously been mistaken for a library primitive.

\begin{longtable}{@{}p{0.52\textwidth}p{0.42\textwidth}@{}}
\toprule
\textbf{Not in the base} & \textbf{Derived in} \\
\midrule
\endhead
Signed Riesz representation (\ref{base:measure} gives positive functionals
only; \texttt{RealRMK.rieszMeasure} returns a \texttt{Measure}, not a
\texttt{SignedMeasure})
  & \S\ref{sec:signed-riesz} \\
Weak-$*$ compactness stated for signed or vector Radon measures
  & \S\ref{sec:signed-compactness} \\
Sets of finite perimeter, reduced boundary, several-variable BV
(mathlib's \code{Analysis/BoundedVariation.lean} is one-dimensional)
  & Ch.~\ref{ch:bv}--\ref{ch:degiorgi} \\
The \emph{lower-dimensional} (rectifiable) area formula; \ref{base:calculus}
gives the equidimensional one only
  & \S\ref{sec:area-formula} \\
De Giorgi's structure theorem; \ref{base:besicovitch} is not a substitute
  & Ch.~\ref{ch:degiorgi} \\
Coarea for BV and for $C^1$ functions with a weight
  & \S\ref{sec:coarea} \\
Any Sobolev, Poincar\'e, Rellich or isoperimetric inequality
  & Ch.~\ref{ch:sobolev} \\
Elliptic estimates: Campanato, Schauder, maximum and Hopf principles,
boundary regularity
  & Ch.~\ref{ch:elliptic} \\
Sard's theorem in any form
  & Ch.~\ref{ch:sard} \\
Differentiability of ODE solutions in the initial point; a global flow
generated by a vector field (\texttt{Dynamics/Flow.lean} is an abstract
continuous monoid action, never built from an ODE); the flow-box theorem
  & \S\ref{sec:flow} \\
Transversality in any form; tubular neighbourhoods
  & \S\ref{sec:transversality}, \S\ref{sec:tubular} \\
Morse theory; Gauss--Bonnet; the degree of a smooth map; classification of
compact surfaces
  & \S\ref{sec:shape}--\S\ref{sec:total-curvature} (degree theory is
    \emph{avoided}, not derived --- Formalisation
    note~\ref{fnote:not-a-degree}) \\
Any regularity theorem for perimeter (quasi)minimisers
  & Ch.~\ref{ch:eps-reg} \\
Capacitary monotonicity (Agostiniani--Mazzieri)
  & Ch.~\ref{ch:appK} \\
\bottomrule
\end{longtable}

\begin{fnote}[The ODE trap]\label{fnote:ode-trap}
The mathlib ODE results are named by the shape of their conclusion, not by
the words ``dependence'' or ``flow'', and the local flow they produce is a
bare function rather than a \texttt{Flow}.  Keyword searches miss them.
Grep for conclusion shapes, or read
\code{Mathlib/Analysis/ODE/ExistUnique.lean} directly.  In particular
\ref{base:ode} \emph{does} give Lipschitz dependence on the initial point and
joint continuity, and \emph{does not} give differentiability in the initial
point.
\end{fnote}

\begin{fnote}[Two false friends]\label{fnote:false-friends}
Do not use \ref{base:calculus}'s equidimensional Jacobian theorem where the
rectifiable area formula (\S\ref{sec:area-formula}) is required, and do not
use \ref{base:besicovitch}'s differentiation theorem where De Giorgi's
structure theorem (Ch.~\ref{ch:degiorgi}) is required.  Both substitutions
typecheck as English and neither is valid.
\end{fnote}

